%杨舒云的实验报告编辑界面，使用了Huanyu Shi，2019级的模板，杨舒云在此拜谢ORZ!

%!TEX program = xelatex
\documentclass[dvipsnames, svgnames,a4paper,11pt]{article}
\input{YsySettings} 
\usepackage{lipsum}

\begin{document}
	


%实验报告封面	
	\begin{table}
		\renewcommand\arraystretch{1.7}
		\begin{tabularx}{\textwidth}{
				|X|X|X|X
				|X|X|X|X|}
			\hline
			\multicolumn{2}{|c|}{预习报告}&\multicolumn{2}{|c|}{实验记录}&\multicolumn{2}{|c|}{分析讨论}&\multicolumn{2}{|c|}{总成绩}\\
			\hline
			& & & & & & & \\
			\hline
		\end{tabularx}
	\end{table}
		
	\begin{table}
		\renewcommand\arraystretch{1.7}
		\begin{tabularx}{\textwidth}{|X|X|X|X|}
			\hline
			专业： & 物理学 &年级： & 2022级\\
			\hline
			姓名： & 杨舒云  & 学号： & 22344020\\
			\hline
			实验时间： & 2023/9/7 & 教师签名： & \\
			\hline
		\end{tabularx}
	\end{table}

	\begin{center}
		\LARGE BC2 \quad 薄透镜焦距的测量
	\end{center}
	
	\textbf{【实验报告注意事项】}
	\begin{enumerate}
		\item 实验报告由三部分组成：
		\begin{enumerate}
			\item 预习报告：（提前一周）认真研读\underline{\textbf{实验讲义}}，弄清实验原理；实验所需的仪器设备、用具及其使用（强烈建议到实验室预习），完成课前预习思考题；了解实验需要测量的物理量，并根据要求提前准备实验记录表格（第一循环实验已由教师提供模板，可以打印）。预习成绩低于10分（共20分）者不能做实验。
			\item 实验记录：认真、客观记录实验条件、实验过程中的现象以及数据。实验记录请用珠笔或者钢笔书写并签名（\textcolor{red}{\textbf{用铅笔记录的被认为无效}}）。\textcolor{red}{\textbf{保持原始记录，包括写错删除部分，如因误记需要修改记录，必须按规范修改。}}（不得输入电脑打印，但可扫描手记后打印扫描件）；离开前请实验教师检查记录并签名。
			\item 分析讨论：处理实验原始数据（学习仪器使用类型的实验除外），对数据的可靠性和合理性进行分析；按规范呈现数据和结果（图、表），包括数据、图表按顺序编号及其引用；分析物理现象（含回答实验思考题，写出问题思考过程，必要时按规范引用数据）；最后得出结论。
		\end{enumerate}
		\textbf{实验报告就是将预习报告、实验记录、和数据处理与分析合起来，加上本页封面。}
		\item 每次完成实验后的一周内交\textbf{实验报告}（特殊情况不能超过两周）。
		\item 除实验记录外，实验报告其他部分建议双面打印。
	\end{enumerate}
	
	\textbf{【实验安全注意事项】}	
	\begin{enumerate}
		\item 实验中需带手套，尽量避免用手直接接触镜片的光学面。若不小心触摸了光学表面，需尽快用镜头纸或擦镜布擦拭干净。	
		\item 安装镜片时需在光学平台上尽量靠近台面的高度操作，以免失手跌落摔碎镜片。
		\item 实验平台配件所用固定螺钉需拧紧，以免镜架晃动；但不可过紧，以免损坏。
		\item 实验前需按仪器清单检查光学元件是否齐全，实验结束后按照顺序放回元件盒。
	\end{enumerate}

%	\textbf{【特别鸣谢】}	
	
%	感谢2019级学长石寰宇为本实验报告提供\LaTeX 模板


 
	\clearpage
	\tableofcontents
	\clearpage



%预习报告	
	\setcounter{section}{0}
	\section{BC2 薄透镜焦距的测量 \quad\heiti 预习报告}
	
	\subsection{实验目的}
	\begin{enumerate}
		\item 熟悉光具座及其附件的使用。
		\item 通过实验进一步理解薄透镜的成像规律。
		\item 掌握利用自准直方法测量薄透镜焦距。
		\item 掌握和理解光学系统共轴等高调节的方法。
	\end{enumerate}
	
	\subsection{仪器用具}
	\begin{table}[htbp]
		\centering
		\renewcommand\arraystretch{1.6}
		% \setlength{\tabcolsep}{10mm}
		\begin{tabular}{p{0.05\textwidth}|p{0.20\textwidth}|p{0.05\textwidth}|p{0.5\textwidth}}
			\hline
			编号& 仪器用具名称 & 数量 &  主要参数（型号，测量范围，测量精度等） \\
			\hline
			1& 透镜LC &1 & \\
			
			2& 白光光源 &1 & GY-6 型，亮度可调，即溴钨灯 \\
			
			3& 白屏 & 1 & SZ-13 \\
			
			4& 物屏 & 1 & SZ-14 \\
			
			5& 透镜L & 1 & f=150mm 等 \\
			
			6& 透镜架 & 1 & \\
			
			7& 平面镜 & 1 & Φ36*4 \\
			
			8& 平移底座 & 4 & \\
			\hline
		\end{tabular}
	\end{table}
	
	\subsection{原理概述}
	薄透镜是一种光学元件，其特点是其厚度远小于焦距和球面曲率半径，因此常被用于各种光学仪器中。通常，透镜制作材料为光学玻璃，且具有两个折射面。根据折射面的形状，透镜可分为球面透镜和非球面透镜两种。
	
	对于球面透镜，其折射面可以是凸球面、平凸或平面，从而形成双凸、平凸和凸四三种基本形状的透镜。这些透镜能够将平行入射光线聚焦到一点，因此它们具有正的焦距值，被称为会聚透镜或正透镜。
	
	另一方面，非球面透镜的折射面不是球面，包括双四、平四和负弯月(四凸)三种基本形状的透镜。这些透镜会使平行入射光线发散，导致负的焦距值，因此它们被称为发散透镜或负透镜。
	
	1.薄透镜成像公式：
	
	薄透镜成像公式用于描述薄透镜的成像原理。这个公式基于图1，适用于薄透镜，其中我们假设透镜非常薄以至于近似平面。
	
	公式为：
	
	\[
	\frac{1}{f} = \frac{1}{v} + \frac{1}{u}
	\]
	
	其中：
	
	\(f\) 代表薄透镜的焦距。
	
	\(v\) 代表像的位置，正值表示实际像，负值表示虚拟像。
	
	\(u\) 代表物体的位置，正值表示物体在透镜的一侧，负值表示物体在透镜的另一侧。
	
	假设物体在透镜左边，光线由左向右传播，可对x和x'的正负号作如下约定：当物点A与透镜光心距离大于物方焦距f，即A点在F点左侧，则x>0,反之x<0；
	当像点A'与透镜光心距离大于像方焦距f'，即A'点在F'右侧，则x'>0，反之x'<0.
	
	\begin{figure}[htbp]
		\centering
		\includegraphics[width=\textwidth]{BC2Graph1.jpg}
		\caption{薄凸透镜成像原理示意图}
	\end{figure}
	
	2.平行光法测透镜焦距：
	
	根据焦距的定义，测量透镜焦距的最直接方法是使用平行光垂直照射透镜表面，然后改变观察屏与透镜之间的距离。通过找到使光线汇聚至光斑直径最小的位置，可以确定该距离即为透镜的焦距。在实际实验中，可以使用平行光源，如太阳光、距离透镜焦距大于100倍焦距的光源，或者由两个共焦透镜组合产生的光线，来实现平行光的照射。
	
	3.自准法测薄透镜焦距：
	
	自准法是一种用于测量薄透镜焦距的方法，通过观察物体和像的位置，可以确定焦距。以下是该方法的步骤和示意图：
	
	\begin{figure}[htbp]
		\centering
		\includegraphics[width=\textwidth]{BC2Graph2.jpg}
		\caption{自准法测薄透镜焦距实验原理示意图}
	\end{figure}
	
	透镜的像方焦距等于物方焦距的情况发生在自准法测薄透镜焦距实验中。当物体 AB 位于透镜的物方焦平面上时，反射光所形成的像 A'B' 是等大的、倒立的实像，并且 A、A'、F 三点重合在一起。这时，物体与像之间的距离 \(u\) 等于透镜的物方焦距，同时像方焦距 \(v\) 也等于透镜的物方焦距，因为根据薄透镜公式 \(\frac{1}{f} = \frac{1}{v} + \frac{1}{u}\)，当 \(v = u = f\) 时，两者相等。
	
	自准法测薄透镜焦距的光路如图2所示，在透镜的像方放置平面反射镜与光轴垂直，实验中不断改变物AB与透镜之间的距离u，当物AB处在透镜的物方焦平面上时，反射光所成的像是与物等大的倒立实像A'B'，且A、A'、F三点重合分别作出了轴线上的物点和非轴线上物点的成像光路。成像最清晰时的物距u就是物方焦距f. 像方焦距f'的测量方法与物方焦距相同。
	
	步骤：
	
	放置物体（例如箭头）在薄透镜的一侧（\(u\)）。
	
	调整薄透镜的位置，使屏幕上或墙壁上的像清晰可见。
	
	记录物体到薄透镜的距离（\(u\)）和像到薄透镜的距离（\(v\)）。
	
	4.物像放大法测透镜焦距
	实验中，我们使用一个物体 \(O\) 和一个观察屏 \(P\)，并将透镜 \(L\) 放置在物体 \(O\) 和观察屏 \(P\) 之间，以便通过透镜形成像 \(I\)。
	
	放大倍数（\(M\)）：放大倍数表示像 \(I\) 的高度与物体 \(O\) 的高度之比，通常用 \(M\) 表示。它可以通过以下公式计算：
	
	\[
	M = \frac{{\text{像的高度}}}{{\text{物体的高度}}}
	\]
	
	透镜公式：透镜公式用于描述透镜的成像规律，它表示放大倍数 \(M\) 与物体到透镜的距离 \(d_o\)、像到透镜的距离 \(d_i\) 以及焦距 \(f\) 之间的关系：
	
	\[
	M = \frac{{-d_i}}{{d_o}}
	\]
	
	焦距测量：通过测量物体 \(O\) 和像 \(I\) 的高度，以及物体到透镜的距离 \(d_o\) 和像到透镜的距离 \(d_i\)，可以使用透镜公式来计算透镜的焦距 \(f\)。
	
	实验原理的核心思想是，通过测量物体和像的高度以及物体到透镜和像到透镜的距离，可以计算出透镜的焦距，从而实现焦距的测量。
	
	\subsection{实验前思考题}
	\begin{question}
		什么是透镜的“焦点”、“焦距”和“焦平面”？会聚透镜和发散透镜的区别？什么情况下透镜的像方焦距等于透镜的物方焦距？
	\end{question}
	（1）焦点、焦距和焦平面：
	
	焦点（Focal Point）：透镜的焦点是指光线在通过透镜后相交或发散的位置。会聚透镜有一个正焦点，光线汇聚在其前方，而发散透镜有一个负焦点，光线看起来似乎来自于其后方。
	
	焦距（Focal Length）：焦距是指从透镜到其焦点的距离。对于会聚透镜，焦距为正数，对于发散透镜，焦距为负数。
	
	焦平面（Focal Plane）：焦平面是与透镜焦点垂直的平面，光线通过透镜后会在焦平面上交汇或发散。
	
	（2）会聚透镜和发散透镜的区别：
	
	会聚透镜（正透镜）：会聚透镜将平行光线折射后使其汇聚到一个点，这个透镜有正的焦距。它通常用于放大或聚焦光线，如凸透镜。
	
	发散透镜（负透镜）：发散透镜使平行光线折射后发散开，看起来似乎来自透镜后方，有负的焦距。它通常用于减小或分散光线，如凹透镜。
	
	（3）透镜的像方焦距等于物方焦距的情况：
	
	在物体和像所在的两个不同介质材料中，双方的焦距通常是不相等的。
	
	透镜的焦距与折射率有关，具体取决于透镜所在的介质以及物体和像所在的介质。因此，在通常情况下，如果物体和像分别位于两种不同的介质中，它们的折射率不同，导致透镜的物方焦距和像方焦距也不相同。透镜的焦距是根据两个介质之间的折射率差异来计算的。
	
	只有当物体和像都位于相同的介质中，或者两个介质的折射率相等时，物方焦距和像方焦距才会相等。这种情况下，透镜的焦距在这两个方向上是相同的。
	
	进一步来说，物体在一个介质中，而像在另一个介质中时，透镜的物方焦距和像方焦距通常不相等。物体和像都在同一介质中时，透镜的物方焦距和像方焦距可以相等。在这两种情况下，透镜焦距的差异取决于两个介质的折射率。根据透镜公式，焦距与物方焦距和像方焦距之间的关系如下：
	
	\[
	\frac{1}{f} = (n_2 - n_1) \left( \frac{1}{R_1} - \frac{1}{R_2} \right)
	\]
	
	其中，\(f\) 是透镜焦距，\(n_1\) 和 \(n_2\) 分别是物体和像所在介质的折射率，\(R_1\) 和 \(R_2\) 分别是透镜两侧的曲率半径。根据这个公式，不同的折射率将导致不同的焦距。
	
	总之，在不同介质中，通常情况下，透镜的物方焦距和像方焦距不相等。但如果物体和像都在同一介质中，它们可能会相等。
	
	\begin{question}
		本实验的误差主要来源，如何减小误差？
	\end{question}
	
	（1）准直光线的调整误差：在实验中，确保准直光线的准确调整是关键。误差可能由于准直光线的不精确调整而引入，导致物体和像的位置测量不准确。
	
	减小误差方法：使用高质量的光源和适当的准直装置，确保光线是平行的。使用准确的光源位置和准直器来调整准直光线。
	
	（2）反射镜的位置稳定性：平面反射镜的位置应保持稳定，以确保反射光线正确返回。反射镜位置的微小变化可能导致像的位置变化。
	
	减小误差方法：使用稳定的支架和固定反射镜，以防止其移动或摆动。确保反射镜的位置在实验过程中保持不变。
	
	（3）读数误差：读取物体和像的位置可能会引入误差。这些误差可能来自于尺度的不准确度或读数时的视觉偏差。
	
	减小误差方法：使用高精度的测量工具，如数字卡尺，以提高读数的准确性。进行多次测量并取平均值，以减小随机误差。
	
	（4）物体大小：物体的实际大小可能与理想的点物体假设不符，这可能影响焦点的精确位置。
	
	减小误差方法：选择尽量小的物体或使用点光源，以减小物体大小对实验的影响。
	
	（5）折射率的变化：光线在透镜材料中的传播受折射率影响，而折射率可能随温度或波长而变化。
	
	减小误差方法：在实验期间保持透镜和环境温度稳定，并使用单色光源以减小波长变化。
	
	综合来看，误差来自于：物体位置调整误差——调整物体的位置时，可能存在微小的误差，导致焦点位置不准确。像位置测量误差——测量像的位置时，也可能存在位置测量误差，特别是当光斑边缘不明显时。环境条件变化——温度和气压的变化可以影响光线的折射行为，因此需要在实验中控制环境条件。减小误差的办法有：使用精确的位置测量工具——使用高精度的测量工具来测量物体和像的位置，以减小位置测量误差。使用辅助设备——使用适当的辅助设备，如定位平台或显微镜，来帮助精确地调整物体和测量像的位置。多次测量和平均值——进行多次测量并计算平均值，以减小由于随机误差引起的不确定性。控制环境条件——在实验中控制温度和气压等环境条件，以减小环境因素对测量的影响。注意反射——在自准焦法中，要注意是否有反射光线的干扰，确保只测量主要的聚焦光线。
	
	\begin{question}
		采用三条基本光线绘图的方法画出自准法测量薄透镜焦距的实验原理光路图。
	\end{question}
	此问题相关文字说明已在实验原理部分有了详细记载，以下只提供示意图（图3）。
	
	\begin{figure}[htbp]
		\centering
		\includegraphics[width=\textwidth]{BC2Graph2.jpg}
		\caption{思考题之自准法测量薄透镜焦距的实验原理光路图}
	\end{figure}



%实验记录	
	\clearpage
	\begin{table}
		\renewcommand\arraystretch{1.7}
		\centering
		\begin{tabularx}{\textwidth}{|X|X|X|X|}
			\hline
			专业： & 物理学 & 年级： & 2022级 \\
			\hline
			姓名： & 杨舒云 & 学号： & 22344020\\
			\hline
			室温： & 26℃ & 实验地点： & 实验室512 \\
			\hline
			学生签名：& 杨舒云 & 评分： &\\
			\hline
			实验时间：& 2023/9/7 & 教师签名：&\\
			\hline
		\end{tabularx}
	\end{table}
	
	\section{BC2 薄透镜焦距的测量 \quad\heiti 实验记录}
	\subsection{实验内容、步骤与结果}
	\subsubsection{简单测量会聚透镜的焦距}
	\begin{enumerate}
		\item 摆好光路图，调节光学元件等高共轴
		\item 以溴钨灯白光源为物，放置在尽可能远处，通过透镜形成一个实焦点
		\item 测量记录透镜物方焦距7次；反转透镜，再次测量记录透镜像方焦距		
	\end{enumerate}
	
	实验数据如下所示
	
	\begin{figure}[htbp]
		\centering
		\includegraphics[width=\textwidth]{BC2Graph3.png}
		\caption{简单测量会聚透镜的焦距的测量结果}
	\end{figure}
	
	事实上，从上表可以发现测量结果与实验仪器上给定的参考值有很大的偏差，个人认为是因为溴钨灯即使放在较远处也并不能视为平行光源，因此并不能采用直接读出焦距的方法。所以，在采用成像公式修正后，得到了较好的结果，详见\textbf{实验数据分析}部分。
	
	\subsubsection{自准法测会聚透镜的焦距}
	\begin{enumerate}
		\item 摆好光路，将物屏摆在光源前；调节各元件等高共轴。
		\item 移动透镜，简单描述如何判断物屏所在位置即透镜焦距？并测量记录物屏和透镜的位置A1，A2各7次。反转透镜，做必要的调节，再测量记录物屏和透镜的位B1，B2各7次。
		\item 透镜焦距确定之后，改变平面镜与透镜之间的距离，描述物屏上观察的现象。
	\end{enumerate}
	
	实验数据如下所示：
	
	\begin{figure}[htbp]
		\centering
		\includegraphics[width=\textwidth]{BC2Graph4.jpg}
		\caption{自准法测会聚透镜的焦距的测量结果}
	\end{figure}
	注：\(f=A_{2}-A_{1},f^{'}=B_{2}-B_{1}\)
		
	在确定焦距后，只移动平面镜位置，可观察到在超过一定范围后物屏上所成的倒立等大的像开始缩小，并随平面镜距离增加而像的大小不断减小。关于这个现象发生的原理的猜想，详见\textbf{思考题3.2}。    
	
	\subsubsection{物像放大法测透镜的焦距}
	这个实验的内容与步骤是由学生自己设计的，由于实验设备的原因（详见\textbf{问题记录}部分），这个实验是由我（杨舒云）与戴鹏辉（22344016）在他的实验台上共同完成的，我选择了我们得到的数据中的以下几组并得到结果。
	
	实验步骤如下所示：
	
	\begin{figure}[htbp]
		\centering
		\includegraphics[width=\textwidth]{BC2Graph5.jpg}
		\caption{物像放大法测透镜的焦距实验装置图}
	\end{figure}
	
	首先按照图6装好实验仪器，其中，为了更好地控制透镜与微尺板间距离，我们将3和5替换为了如图7所示的节点架。
	
	\begin{figure}[htbp]
		\centering
		\includegraphics[width=\textwidth]{BC2Graph6.jpg}
		\caption{节点架示意图}
	\end{figure}
	
	接下来，点亮光源，先将M、Le、ME紧靠在一起，然后固定Le，让M、ME沿光轴逐渐远离Le，直至在ME中观察到微尺的放大像，再微调ME的位置，使像清晰且与ME微尺板刻线无视差。
	
	测出1/10mm微尺的像中相邻两刻线间的距离h，按实验原理中所述式子计算放大率m，读出M与ME的在光具座上的相对位置a和b。
	
	重复上述步骤得到一系列数据，按公式 \( f=\frac{\left | \left ( b_{2}-b_{1} \right ) -\left ( a_{2}-a_{1} \right )  \right |}{\left | m_{2}-m_{1} \right | } \) 计算出焦距。
	
	实验数据如下所示：
	\begin{figure}[htbp]
		\centering
		\includegraphics[width=\textwidth]{BC2Graph7.jpg}
		\caption{物像放大法测透镜的焦距的测量结果}
	\end{figure}
	
	从数据表可以看出，4、5组数据明显超出3$\sigma$，应该被舍去。
	
	\clearpage
	\subsection{原始数据记录}
	实验记录本上的原始数据、整理后的实验台、预习报告签名如下所示：
	
	\begin{figure}[htbp]
		\centering
		\includegraphics[width=\textwidth]{BC2Graph8-1.jpg}
		\caption{原始数据1}
	\end{figure}
	
	\begin{figure}[htbp]
		\centering
		\includegraphics[width=\textwidth]{BC2Graph8-2.jpg}
		\caption{原始数据2}
	\end{figure}
	
	\begin{figure}[htbp]
		\centering
		\includegraphics[width=\textwidth]{BC2Graph8-3.jpg}
		\caption{原始数据3}
	\end{figure}
	
	\begin{figure}[htbp]
		\centering
		\includegraphics[width=\textwidth]{BC2Graph8-4.jpg}
		\caption{实验台整理}
	\end{figure}
	
	\begin{figure}[htbp]
		\centering
		\includegraphics[width=\textwidth]{BC2Graph8-5.jpg}
		\caption{预习报告签字}
	\end{figure}
	
	\clearpage
	\subsection{实验过程中遇到的问题记录}
	\begin{enumerate}
		\item 如2.1.1中所述，按照实验步骤所得到的测量结果与参考值偏差较大，我询问了其他同学，他们的测量也均偏差较大。
		\item 在2.1.3中我曾谈到了我自己的实验台的装置出现了问题导致我无法正常测量，具体情况是：首先，我使用白屏检测了微尺板是否能正常成像，结果是否定的；接下来，我更换为了其他同学的微尺板，能正常成像，但用读数显微镜仍无法观测到，在这里我猜测是读数显微镜出了问题。为了不影响实验进行，我选择了在他人的操作台上完成实验。
	\end{enumerate}
	


%分析与讨论	
	\clearpage
	\begin{table}
		\renewcommand\arraystretch{1.7}
		\begin{tabularx}{\textwidth}{|X|X|X|X|}
			\hline
			专业：& 物理学 &年级：& 2019级\\
			\hline
			姓名： & 杨舒云 & 学号：& 22344020\\
			\hline
			日期：& 2023/9/7 & 评分： &\\
			\hline
		\end{tabularx}
	\end{table}
	
	\section{BC2 薄透镜焦距的测量 \quad\heiti 分析与讨论}
	\subsection{实验数据分析}
	\begin{enumerate}
		\item 在2.1.1中我已经简要给出了简单测量法测得的会聚透镜的物方焦距和像方焦距和相应的误差分析，详见图4；由于误差过大，我并不认可该方法下的结果，因此我在这里只讨论测量值与理论值的差别的原因，并重点分析修正后的误差。
		
		首先，我认为用点光源的模型去考虑比用平行光的模型去考虑更合理，这正是测量值与参考值偏差较大的原因。因此，考虑点光源模型，采用 \(f=\frac{1}{\frac{1}{u}+\frac{1}{v}}\) 计算得到焦距，其中\(v\)为像距，而\(u\)用光源到透镜的距离来近似，计算结果如下：
		
		\begin{figure}[htbp]
			\centering
			\includegraphics[width=\textwidth]{BC2Graph9.jpg}
			\caption{简单测量会聚透镜的焦距的修正结果}
		\end{figure}
		
		可以看出，相对误差更小，也就是说点光源模型更加符合这个实验。
		
		接下来进一步进行误差分析，考虑各测量列的A类不确定度并计算，然后B类不确定度作如下考虑——光具座的最小分度值为1 mm，故取$u_{b}=\frac{1 mm}{\sqrt{3}}=0.5774mm $，最终得到：
		\begin{equation*}
			f=\bar{f}\pm U_{f}\approx( 14.92 \pm  2.31 )cm
		\end{equation*} 
		其中展伸不确定度$U_{f}={k}\times{u_{f}}\approx 2.31 $ 是由合成标准不确定度$u_{f}\approx 0.98 $ 和$k\approx2.81$确定的，k是依据置信概率$P=99.73\%$和自由度$\nu\approx7.44$查t分布表得到的。
		
		\begin{equation*}
			f^{'}=\bar{f^{'}}\pm U_{f^{'}}\approx( 14.70 \pm  1.24 )cm
		\end{equation*} 
		其中展伸不确定度$U_{f^{'}}={k}\times{u_{f^{'}}}\approx 1.24 $ 是由合成标准不确定度$u_{f^{'}}\approx 0.81 $ 和$k\approx1.37$确定的，k是依据置信概率$P=99.73\%$和自由度$\nu\approx4.57$查t分布表得到的。
		
		\item 自准法测得的会聚透镜的物方焦距、像方焦距和相应的误差分析如下所示：
		
		将均值作为测量值，将实验标准差作为A类不确定，B类不确定度依旧作如下考虑——具座的最小分度值为1 mm，故取$u_{b}=\frac{1 mm}{\sqrt{3}}=0.5774mm $，最终得到：
		\begin{equation*}
			f=\bar{f}\pm U_{f}=\bar{f}\pm\sqrt{u_{A}^{2}+u_{B}^{2}}\approx( 15.04 \pm  1.43 )cm
		\end{equation*} 
		其中展伸不确定度$U_{f}={k}\times{u_{f}}\approx 1.43 $ 是由合成标准不确定度$u_{f}\approx 0.58 $ 和$k\approx4.30$确定的，k是依据置信概率$P=95\%$和自由度$\nu\approx2.74$查t分布表得到的。
		
		\begin{equation*}
			f^{'}=\bar{f^{'}}\pm U_{f^{'}}=\bar{f^{'}}\pm\sqrt{u_{A}^{2}+u_{B}^{2}}\approx( 14.92 \pm  1.44 )cm
		\end{equation*} 
		其中展伸不确定度$U_{f^{'}}={f^{'}}\times{u_{f^{'}}}\approx 1.44 $ 是由合成标准不确定度$u_{f^{'}}\approx 0.58 $ 和$k\approx4.30$确定的，k是依据置信概率$P=95\%$和自由度$\nu\approx2.74$查t分布表得到的。
		
		以上数据利用到了数据表（图5）中数据。
		
		\item 由于我个人认为在完成这个实验时有很多不完善的地方，由此得到的数据偏差会很大，因此只作如下最基本的误差分析。首先在略去数据表（图8）3$\sigma$以外的数据以后，由1、2和3组得到以下结果：
		
		均值为$\bar{f}\approx3.23cm$，实验标准差为$U_{f}\approx0.03cm$，相对误差$\eta_{f}\approx11.52\%$。
		
		未测量$f_{\prime}$的数据。
	\end{enumerate}
		
	\subsection{实验后思考题}
	\begin{question}
		用自准法测量薄透镜焦距时，为什么需要将透镜旋转180度再测量？
	\end{question}
	首先，在思考题1.1中谈到了物方焦距与像方焦距会存在一定的偏差，因此，测量物方焦距后，将透镜旋转180度可再次测量像方焦距。
	
	其次，用自准法测量薄透镜焦距时，若要得到一个焦距的值，可通过这个方法消除透镜光心偏离滑座位置指针而产生的误差。因为透镜的光心位置可能不与滑座的指针位置重合，所以如果只进行单侧测量，就会引入一定的误差。但是如果将透镜旋转180度后再测量一次，那么这种误差就会相互抵消，从而提高测量的准确性。这种方法叫做两次平均法，是一种常用的减小系统误差的技术。
	
	具体来说，在自准法测量薄透镜焦距时，将透镜旋转180度再测量是为了检验透镜的两个折射面是否是相同的，以确保测量的准确性。这是因为透镜的两个折射面可能存在微小的制造差异或不完全对称，这些差异可能会导致焦距测量的误差。通过将透镜旋转180度，实际上改变了光线通过透镜的路径。如果透镜的两个折射面是完全对称的，那么无论透镜是正面朝上还是下，焦点位置应该是相同的。但如果存在微小差异，透镜旋转后焦点位置可能会略有变化。通过比较两次测量的结果，您可以检测出任何差异，然后可以取两次测量的平均值，从而减小由于透镜不对称性引起的误差。这种方法有助于提高焦距测量的准确性，特别是当需要高精度测量透镜焦距时。
	
	\begin{question}
		自准直法中平面镜离透镜远近对测量是否有影响？
	\end{question}
	理论上，根据原理图（图2），平面镜的远近对测量没有影响。通常情况下，为了获得清晰的焦点和准确的测量结果，平面镜应该相对于透镜的焦点位置足够远，以确保只观察到主要的聚焦光线，而不受多次反射或越过反射镜的影响。
	
	但在实际上，由于平面镜和透镜的大小限制，平面镜离得太远会导致部分透过透镜的光线无法被平面镜反射而回，导致本该等大的像实际上会缩小。
	\begin{figure}[htbp]
		\centering
		\includegraphics[width=\textwidth]{BC2Graph10.jpg}
		\caption{我们所使用的平面镜}
	\end{figure}
	除此之外，进一步考虑，距离过小——如果平面镜离透镜太近，光线可能会发生多次反射，导致焦点不明确，从而使测量变得更加复杂和不准确。距离过大——如果平面镜离透镜太远，光线可能会越过反射镜，使反射光线不容易观察到，这也会导致焦点的定位困难。
	
	\begin{question}
		光路的共轴调节要注意哪些？
	\end{question}
	光路的共轴调节是指使各个光学元件的光心在同一条直线上，从而保证光束的传输和成像的质量。光路的共轴调节要注意以下几点：
	
	共轴调节一般分为粗调和细调两个步骤，粗调时要使各个光学元件的中心大致在一条直线上，且与光具座的导轴平行，细调时要用小孔或光阑来检验是否共轴。
	
	共轴调节时要注意光学元件的方向和位置，如透镜要使平面朝向焦点，反射镜要垂直于光轴，分束镜要与入射光成45度角等。
	
	共轴调节时要注意光束的质量和准直度，如使用扩束镜或针孔滤波来增加光束直径，使用剪切干涉仪来验证光束是否准直等。
	
	对于多透镜系统，可以先将每个透镜单独调节到共轴，然后再将它们组合起来，或者先将它们组合成一个整体，然后再进行共轴调节。对于复杂的透镜系统，可以用激光散斑法来确定透镜的焦平面位置。
	
	综上所述，以下是一些要注意的要点：确保光源的位置稳定，通常需要将光源放置在远离实验区域，以避免干扰光路。确保透镜的位置准确，透镜的中心应与光轴对齐。调整物体的位置，以确保物体位于透镜的物方焦平面上。防止光路中的元件或支架发生移动或振动，这可能会影响焦距的测量。	使用合适的测量工具（如尺子或标尺）来测量物体到透镜的距离，以获得准确的数据。小心处理透镜和其他光学元件，以避免划伤或损坏它们。如果实验需要在不同位置观察或测量物体，确保光学元件的位置可以通过调焦装置进行精确调整。
	
	\begin{question}
		用物像放大倍率测量透镜焦距的实验中，测微目镜的物面在什么位置？作图说明。
	\end{question}
	用物像放大倍率测量透镜焦距的实验中，测微目镜的物面应该在待测目镜的像方焦平面上，也就是说，测微目镜的物距等于待测目镜的像方焦距。这样可以使待测目镜的像与测微目镜的标尺无视差，从而准确地测量像的放大倍率。换句话说，在用物像放大倍率测量透镜焦距的实验中，测微目镜的物面通常位于透镜的物方焦点上。这意味着测微目镜中的物体应该被放置在透镜的物方焦点处，以便能够观察到透镜形成的实际像。在这个位置，测微目镜将能够以放大倍率观察到焦点附近的物体，从而用于测量透镜的焦距。
	\begin{figure}[htbp]
		\centering
		\includegraphics[width=\textwidth]{BC2Graph11.jpg}
		\caption{物像放大倍率测量透镜焦距光路图}
	\end{figure}
	
	\section{BC2 薄透镜焦距的测量 \quad\heiti 参考文献}
	基础物理实验/沈韩主编.——北京：科学出版社，2015.2 ISBN：978-7-03-043311-4 P236~241
	
	
	
\end{document}